\documentclass[10pt,oneside]{article}
\input{chapters-tex/shared/book-metadata}
\usepackage[T1]{fontenc}
\usepackage[utf8]{inputenc}
\usepackage{newpxtext}
\usepackage[scaled=.94]{sourcesanspro}
\input{chapters-tex/shared/daft-macros}
\let\BookTranspose\transp\let\transp\relax\usepackage{newpxmath}\let\transp\BookTranspose
\input{ursfs.fd}\DeclareFontShape{U}{rsfs}{m}{n}{<5-7>rsfs5 <7-8>rsfs7 <8->rsfs10}{}
\usepackage[a4paper,landscape,margin=15mm]{geometry}
\usepackage{xcolor,hyperref,fancyhdr}
\definecolor{BookAccent}{HTML}{146C78}
\hypersetup{colorlinks=true,urlcolor=BookAccent}
\pagestyle{fancy}\fancyhf{}\fancyfoot[L]{\sffamily\small DAFT --- figure comparison}\fancyfoot[R]{\sffamily\small\thepage}
\setlength{\parindent}{0pt}\setlength{\emergencystretch}{2em}
\begin{document}
{\sffamily\color{BookAccent}\LARGE\bfseries Before and after}\par\medskip
{\large\BookTitle\par}\smallskip
\BookAuthor\quad---\quad\BookEditionDate\par\bigskip
The left panels reproduce the previous English edition. The right panels are regenerated from Python or TikZ and used in the revised book. Figure numbers and captions below are read from the compiled book and its source. The original number is also recorded when needed.\par\bigskip
\clearpage
{\sffamily\color{BookAccent}\Large\bfseries Figure 1.1}\par\medskip
Real and imaginary parts of the first four characters of $\ZZ/12\ZZ$.\par\smallskip
{\small Book page 2; original figure 1.1.}\par\bigskip
\begin{minipage}[t]{.485\linewidth}\centering{\sffamily\bfseries Previous edition}\par\medskip
\includegraphics[width=\linewidth,height=.66\textheight,keepaspectratio]{comparator-figures/before/fig-cyclic-characters.pdf}
\end{minipage}\hfill\begin{minipage}[t]{.485\linewidth}\centering{\sffamily\bfseries Revised edition}\par\medskip
\includegraphics[width=\linewidth,height=.66\textheight,keepaspectratio]{figures/chapter-finite-groups/cyclic-characters/figure.pdf}
\end{minipage}
\clearpage
{\sffamily\color{BookAccent}\Large\bfseries Figure 1.2}\par\medskip
An even bell-shaped function on $\ZZ/17\ZZ$ and its real Fourier coefficients, with centered indices.\par\smallskip
{\small Book page 7; original figure 1.2.}\par\bigskip
\begin{minipage}[t]{.485\linewidth}\centering{\sffamily\bfseries Previous edition}\par\medskip
\includegraphics[width=\linewidth,height=.66\textheight,keepaspectratio]{comparator-figures/before/fig-decomposition-fourier.pdf}
\end{minipage}\hfill\begin{minipage}[t]{.485\linewidth}\centering{\sffamily\bfseries Revised edition}\par\medskip
\includegraphics[width=\linewidth,height=.66\textheight,keepaspectratio]{figures/chapter-finite-groups/fourier-decomposition/figure.pdf}
\end{minipage}
\clearpage
{\sffamily\color{BookAccent}\Large\bfseries Figure 1.3}\par\medskip
Circular self-convolution of $f=\mathbf{1}_{\{5,\ldots,10\}}$ on $\ZZ/16\ZZ$. The maximum occurs at $15$ modulo $16$; the remaining values wrap around at $0$.\par\smallskip
{\small Book page 8; original figure 1.3.}\par\bigskip
\begin{minipage}[t]{.485\linewidth}\centering{\sffamily\bfseries Previous edition}\par\medskip
\includegraphics[width=\linewidth,height=.66\textheight,keepaspectratio]{comparator-figures/before/fig-convolution-intro.pdf}
\end{minipage}\hfill\begin{minipage}[t]{.485\linewidth}\centering{\sffamily\bfseries Revised edition}\par\medskip
\includegraphics[width=\linewidth,height=.66\textheight,keepaspectratio]{figures/chapter-finite-groups/circular-convolution/figure.pdf}
\end{minipage}
\clearpage
{\sffamily\color{BookAccent}\Large\bfseries Figure 1.4}\par\medskip
Random walks on $\ZZ/15\ZZ$ with increment laws $(0,1/2,0,\ldots,0,1/2)$ and $(0.3,0.3,0.2,0,\ldots,0,0.1,0.1)$. Dashed lines show the uniform distribution.\par\smallskip
{\small Book page 12; original figure 1.4.}\par\bigskip
\begin{minipage}[t]{.485\linewidth}\centering{\sffamily\bfseries Previous edition}\par\medskip
\includegraphics[width=\linewidth,height=.66\textheight,keepaspectratio]{comparator-figures/before/fig-marche-aleatoire.pdf}
\end{minipage}\hfill\begin{minipage}[t]{.485\linewidth}\centering{\sffamily\bfseries Revised edition}\par\medskip
\includegraphics[width=\linewidth,height=.66\textheight,keepaspectratio]{figures/chapter-finite-groups/random-walks/figure.pdf}
\end{minipage}
\clearpage
{\sffamily\color{BookAccent}\Large\bfseries Figure 2.1}\par\medskip
The Boolean cube $\FF_2^4$, arranged by Hamming weight. Each edge changes exactly one coordinate.\par\smallskip
{\small Book page 19; original figure 2.1.}\par\bigskip
\begin{minipage}[t]{.485\linewidth}\centering{\sffamily\bfseries Previous edition}\par\medskip
\includegraphics[width=\linewidth,height=.66\textheight,keepaspectratio]{comparator-figures/before/fig-cube-booleen.pdf}
\end{minipage}\hfill\begin{minipage}[t]{.485\linewidth}\centering{\sffamily\bfseries Revised edition}\par\medskip
\includegraphics[width=\linewidth,height=.66\textheight,keepaspectratio]{figures/chapter-duality/boolean-cube/figure.pdf}
\end{minipage}
\clearpage
{\sffamily\color{BookAccent}\Large\bfseries Figure 2.2}\par\medskip
The Walsh matrix $W_{64}$ in natural order. White represents $+1$ and dark teal represents $-1$.\par\smallskip
{\small Book page 20; original figure 2.2.}\par\bigskip
\begin{minipage}[t]{.485\linewidth}\centering{\sffamily\bfseries Previous edition}\par\medskip
\includegraphics[width=\linewidth,height=.66\textheight,keepaspectratio]{comparator-figures/before/fig-matrices-walsh.pdf}
\end{minipage}\hfill\begin{minipage}[t]{.485\linewidth}\centering{\sffamily\bfseries Revised edition}\par\medskip
\includegraphics[width=\linewidth,height=.66\textheight,keepaspectratio]{figures/chapter-duality/walsh-matrix/figure.pdf}
\end{minipage}
\clearpage
{\sffamily\color{BookAccent}\Large\bfseries Figure 2.3}\par\medskip
Lattice-point pairing in the geometric proof of quadratic reciprocity, illustrated with $p=13$, $r=11$, and $a=10$.\par\smallskip
{\small Book page 25; original figure 2.3.}\par\bigskip
\begin{minipage}[t]{.485\linewidth}\centering{\sffamily\bfseries Previous edition}\par\medskip
\includegraphics[width=\linewidth,height=.66\textheight,keepaspectratio]{comparator-figures/before/fig-demo-geom-reciprocite.pdf}
\end{minipage}\hfill\begin{minipage}[t]{.485\linewidth}\centering{\sffamily\bfseries Revised edition}\par\medskip
\includegraphics[width=\linewidth,height=.66\textheight,keepaspectratio]{figures/chapter-duality/geometric-reciprocity/figure.pdf}
\end{minipage}
\clearpage
{\sffamily\color{BookAccent}\Large\bfseries Figure 2.4}\par\medskip
Successive Haar projections onto $F_0$, $F_1$, and $F_2$ (left), and the orthonormal basis functions added at each stage (right).\par\smallskip
{\small Book page 27; original figure 2.4.}\par\bigskip
\begin{minipage}[t]{.485\linewidth}\centering{\sffamily\bfseries Previous edition}\par\medskip
\includegraphics[width=\linewidth,height=.66\textheight,keepaspectratio]{comparator-figures/before/fig-decomposition-base-haar.pdf}
\end{minipage}\hfill\begin{minipage}[t]{.485\linewidth}\centering{\sffamily\bfseries Revised edition}\par\medskip
\includegraphics[width=\linewidth,height=.66\textheight,keepaspectratio]{figures/chapter-duality/haar-projections/figure.pdf}
\end{minipage}
\clearpage
{\sffamily\color{BookAccent}\Large\bfseries Figure 2.5}\par\medskip
The Haar cascade: each stage separates scaling coefficients from details, which are stored from fine to coarse.\par\smallskip
{\small Book page 28; original figure 2.5.}\par\bigskip
\begin{minipage}[t]{.485\linewidth}\centering{\sffamily\bfseries Previous edition}\par\medskip
\includegraphics[width=\linewidth,height=.66\textheight,keepaspectratio]{comparator-figures/before/fig-decomposition-wavelet-haar.pdf}
\end{minipage}\hfill\begin{minipage}[t]{.485\linewidth}\centering{\sffamily\bfseries Revised edition}\par\medskip
\includegraphics[width=\linewidth,height=.66\textheight,keepaspectratio]{figures/chapter-duality/haar-cascade/figure.pdf}
\end{minipage}
\clearpage
{\sffamily\color{BookAccent}\Large\bfseries Figure 2.6}\par\medskip
Two length-$32$ signals and their orthonormal Haar coefficients. Colored separators mark the detail scales; the final coefficient is the scaled mean.\par\smallskip
{\small Book page 28; original figure 2.6.}\par\bigskip
\begin{minipage}[t]{.485\linewidth}\centering{\sffamily\bfseries Previous edition}\par\medskip
\includegraphics[width=\linewidth,height=.66\textheight,keepaspectratio]{comparator-figures/before/fig-transforme-haar.pdf}
\end{minipage}\hfill\begin{minipage}[t]{.485\linewidth}\centering{\sffamily\bfseries Revised edition}\par\medskip
\includegraphics[width=\linewidth,height=.66\textheight,keepaspectratio]{figures/chapter-duality/haar-examples/figure.pdf}
\end{minipage}
\clearpage
{\sffamily\color{BookAccent}\Large\bfseries Figure 2.7}\par\medskip
Walsh functions in natural order and in increasing sequency, with one basis function per row.\par\smallskip
{\small Book page 29; original figure 2.7.}\par\bigskip
\begin{minipage}[t]{.485\linewidth}\centering{\sffamily\bfseries Previous edition}\par\medskip
\includegraphics[width=\linewidth,height=.66\textheight,keepaspectratio]{comparator-figures/before/fig-functions-walsh.pdf}
\end{minipage}\hfill\begin{minipage}[t]{.485\linewidth}\centering{\sffamily\bfseries Revised edition}\par\medskip
\includegraphics[width=\linewidth,height=.66\textheight,keepaspectratio]{figures/chapter-duality/walsh-orderings/figure.pdf}
\end{minipage}
\clearpage
{\sffamily\color{BookAccent}\Large\bfseries Figure 2.8}\par\medskip
The Walsh spectrum of the oscillatory signal in Figure~2.9, in natural and sequency order.\par\smallskip
{\small Book page 29; original figure 2.8.}\par\bigskip
\begin{minipage}[t]{.485\linewidth}\centering{\sffamily\bfseries Previous edition}\par\medskip
\includegraphics[width=\linewidth,height=.66\textheight,keepaspectratio]{comparator-figures/before/fig-spectrum-walsh-1d.pdf}
\end{minipage}\hfill\begin{minipage}[t]{.485\linewidth}\centering{\sffamily\bfseries Revised edition}\par\medskip
\includegraphics[width=\linewidth,height=.66\textheight,keepaspectratio]{figures/chapter-duality/walsh-spectrum/figure.pdf}
\end{minipage}
\clearpage
{\sffamily\color{BookAccent}\Large\bfseries Figure 2.9}\par\medskip
Progressive approximation of a signal by initial segments of the sequency-ordered Walsh basis. Each panel gives the exact retained coefficient count.\par\smallskip
{\small Book page 30; original figure 2.9.}\par\bigskip
\begin{minipage}[t]{.485\linewidth}\centering{\sffamily\bfseries Previous edition}\par\medskip
\includegraphics[width=\linewidth,height=.66\textheight,keepaspectratio]{comparator-figures/before/fig-compression-walsh-1d.pdf}
\end{minipage}\hfill\begin{minipage}[t]{.485\linewidth}\centering{\sffamily\bfseries Revised edition}\par\medskip
\includegraphics[width=\linewidth,height=.66\textheight,keepaspectratio]{figures/chapter-duality/walsh-signal-approximation/figure.pdf}
\end{minipage}
\clearpage
{\sffamily\color{BookAccent}\Large\bfseries Figure 2.10}\par\medskip
The $64$ two-dimensional Walsh functions for $8\times8$ images, ordered by total sequency and then vertical sequency, from left to right and top to bottom.\par\smallskip
{\small Book page 31; original figure 2.10.}\par\bigskip
\begin{minipage}[t]{.485\linewidth}\centering{\sffamily\bfseries Previous edition}\par\medskip
\includegraphics[width=\linewidth,height=.66\textheight,keepaspectratio]{comparator-figures/before/fig-functions-walsh-2d.pdf}
\end{minipage}\hfill\begin{minipage}[t]{.485\linewidth}\centering{\sffamily\bfseries Revised edition}\par\medskip
\includegraphics[width=\linewidth,height=.66\textheight,keepaspectratio]{figures/chapter-duality/walsh-basis-two-dimensions/figure.pdf}
\end{minipage}
\clearpage
{\sffamily\color{BookAccent}\Large\bfseries Figure 2.11}\par\medskip
Walsh approximations of a synthetic letter A, with exact retained coefficient counts. Every panel uses the same grayscale range.\par\smallskip
{\small Book page 32; original figure 2.11.}\par\bigskip
\begin{minipage}[t]{.485\linewidth}\centering{\sffamily\bfseries Previous edition}\par\medskip
\includegraphics[width=\linewidth,height=.66\textheight,keepaspectratio]{comparator-figures/before/fig-compression-walsh-2d.pdf}
\end{minipage}\hfill\begin{minipage}[t]{.485\linewidth}\centering{\sffamily\bfseries Revised edition}\par\medskip
\includegraphics[width=\linewidth,height=.66\textheight,keepaspectratio]{figures/chapter-duality/walsh-image-approximation/figure.pdf}
\end{minipage}
\clearpage
{\sffamily\color{BookAccent}\Large\bfseries Figure 2.12}\par\medskip
Normalized tensor transforms $2^{-n/2}A_\alpha^{\otimes n}$ of a triangular signal, with $n=5$. The middle panel is the Walsh transform; the last reverses the input entries.\par\smallskip
{\small Book page 33; original figure 2.12.}\par\bigskip
\begin{minipage}[t]{.485\linewidth}\centering{\sffamily\bfseries Previous edition}\par\medskip
\includegraphics[width=\linewidth,height=.66\textheight,keepaspectratio]{comparator-figures/before/fig-transfo-walsh-interm.pdf}
\end{minipage}\hfill\begin{minipage}[t]{.485\linewidth}\centering{\sffamily\bfseries Revised edition}\par\medskip
\includegraphics[width=\linewidth,height=.66\textheight,keepaspectratio]{figures/chapter-duality/intermediate-walsh-transform/figure.pdf}
\end{minipage}
\clearpage
{\sffamily\color{BookAccent}\Large\bfseries Figure 3.1}\par\medskip
An FFT butterfly: the same weighted value $wb$ enters the sum and the difference, with $w=\omega_N^{-k}$.\par\smallskip
{\small Book page 36; original figure 3.1.}\par\bigskip
\begin{minipage}[t]{.485\linewidth}\centering{\sffamily\bfseries Previous edition}\par\medskip
\includegraphics[width=\linewidth,height=.66\textheight,keepaspectratio]{comparator-figures/before/fig-fft-butterfly-2points.pdf}
\end{minipage}\hfill\begin{minipage}[t]{.485\linewidth}\centering{\sffamily\bfseries Revised edition}\par\medskip
\includegraphics[width=\linewidth,height=.66\textheight,keepaspectratio]{figures/chapter-discrete-fourier-transform/elementary-butterfly/figure.pdf}
\end{minipage}
\clearpage
{\sffamily\color{BookAccent}\Large\bfseries Figure 3.2}\par\medskip
The combination stage of an eight-point DIT FFT. Two four-point transforms produce $E_k$ and $O_k$; each butterfly combines $E_k$ with $\omega_8^{-k}O_k$.\par\smallskip
{\small Book page 36; original figure 3.2.}\par\bigskip
\begin{minipage}[t]{.485\linewidth}\centering{\sffamily\bfseries Previous edition}\par\medskip
\includegraphics[width=\linewidth,height=.66\textheight,keepaspectratio]{comparator-figures/before/fig-fft-butterfly-8points.pdf}
\end{minipage}\hfill\begin{minipage}[t]{.485\linewidth}\centering{\sffamily\bfseries Revised edition}\par\medskip
\includegraphics[width=\linewidth,height=.66\textheight,keepaspectratio]{figures/chapter-discrete-fourier-transform/eight-point-fft/figure.pdf}
\end{minipage}
\clearpage
{\sffamily\color{BookAccent}\Large\bfseries Figure 3.3}\par\medskip
Bit-reversal permutation matrices for $N=32,64,128,256$. A dot in row $i$ and column $j$ means that the bits of $j$ are those of $i$ in reverse order.\par\smallskip
{\small Book page 38; original figure 3.3.}\par\bigskip
\begin{minipage}[t]{.485\linewidth}\centering{\sffamily\bfseries Previous edition}\par\medskip
\includegraphics[width=\linewidth,height=.66\textheight,keepaspectratio]{comparator-figures/before/fig-rev-bit-matrix.pdf}
\end{minipage}\hfill\begin{minipage}[t]{.485\linewidth}\centering{\sffamily\bfseries Revised edition}\par\medskip
\includegraphics[width=\linewidth,height=.66\textheight,keepaspectratio]{figures/chapter-discrete-fourier-transform/bit-reversal-matrices/figure.pdf}
\end{minipage}
\clearpage
{\sffamily\color{BookAccent}\Large\bfseries Figure 3.4}\par\medskip
Bit reversal on eight entries. The binary labels make the swaps $1\leftrightarrow4$ and $3\leftrightarrow6$ explicit; the other four indices are fixed.\par\smallskip
{\small Book page 38; original figure 3.4.}\par\bigskip
\begin{minipage}[t]{.485\linewidth}\centering{\sffamily\bfseries Previous edition}\par\medskip
\includegraphics[width=\linewidth,height=.66\textheight,keepaspectratio]{comparator-figures/before/fig-inversion-bits.pdf}
\end{minipage}\hfill\begin{minipage}[t]{.485\linewidth}\centering{\sffamily\bfseries Revised edition}\par\medskip
\includegraphics[width=\linewidth,height=.66\textheight,keepaspectratio]{figures/chapter-discrete-fourier-transform/bit-reversal-permutation/figure.pdf}
\end{minipage}
\clearpage
{\sffamily\color{BookAccent}\Large\bfseries Figure 3.5}\par\medskip
Circular and zero-extended linear convolution of a triangular signal with a unit-sum five-tap box filter. The outputs agree away from the endpoints.\par\smallskip
{\small Book page 40; original figure 3.5.}\par\bigskip
\begin{minipage}[t]{.485\linewidth}\centering{\sffamily\bfseries Previous edition}\par\medskip
\includegraphics[width=\linewidth,height=.66\textheight,keepaspectratio]{comparator-figures/before/fig-diff-conv-linear-circular.pdf}
\end{minipage}\hfill\begin{minipage}[t]{.485\linewidth}\centering{\sffamily\bfseries Revised edition}\par\medskip
\includegraphics[width=\linewidth,height=.66\textheight,keepaspectratio]{figures/chapter-discrete-fourier-transform/linear-and-circular-convolution/figure.pdf}
\end{minipage}
\clearpage
{\sffamily\color{BookAccent}\Large\bfseries Figure 3.6}\par\medskip
A two-dimensional checker pattern and its centered Fourier magnitude. The logarithmic color scale reveals the weaker frequencies.\par\smallskip
{\small Book page 41; original figure 3.6.}\par\bigskip
\begin{minipage}[t]{.485\linewidth}\centering{\sffamily\bfseries Previous edition}\par\medskip
\includegraphics[width=\linewidth,height=.66\textheight,keepaspectratio]{comparator-figures/before/fig-tfd-2d.pdf}
\end{minipage}\hfill\begin{minipage}[t]{.485\linewidth}\centering{\sffamily\bfseries Revised edition}\par\medskip
\includegraphics[width=\linewidth,height=.66\textheight,keepaspectratio]{figures/chapter-discrete-fourier-transform/two-dimensional-dft/figure.pdf}
\end{minipage}
\clearpage
{\sffamily\color{BookAccent}\Large\bfseries Figure 3.7}\par\medskip
Circular Gaussian smoothing of a noisy surface. The discrete kernel has unit sum; its center is shifted to index $(0,0)$ before the FFT.\par\smallskip
{\small Book page 42; original figure 3.7.}\par\bigskip
\begin{minipage}[t]{.485\linewidth}\centering{\sffamily\bfseries Previous edition}\par\medskip
\includegraphics[width=\linewidth,height=.66\textheight,keepaspectratio]{comparator-figures/before/fig-convolution-2d.pdf}
\end{minipage}\hfill\begin{minipage}[t]{.485\linewidth}\centering{\sffamily\bfseries Revised edition}\par\medskip
\includegraphics[width=\linewidth,height=.66\textheight,keepaspectratio]{figures/chapter-discrete-fourier-transform/two-dimensional-convolution/figure.pdf}
\end{minipage}
\clearpage
{\sffamily\color{BookAccent}\Large\bfseries Figure 3.8}\par\medskip
The four orthogonal components $P_\lambda f$ of a real signal for the normalized DFT $U$. Their sum is $f$; even and odd components appear on separate rows.\par\smallskip
{\small Book page 43; original figure 3.8.}\par\bigskip
\begin{minipage}[t]{.485\linewidth}\centering{\sffamily\bfseries Previous edition}\par\medskip
\includegraphics[width=\linewidth,height=.66\textheight,keepaspectratio]{comparator-figures/before/fig-square-root-tfd.pdf}
\end{minipage}\hfill\begin{minipage}[t]{.485\linewidth}\centering{\sffamily\bfseries Revised edition}\par\medskip
\includegraphics[width=\linewidth,height=.66\textheight,keepaspectratio]{figures/chapter-discrete-fourier-transform/fourier-eigenspace-projections/figure.pdf}
\end{minipage}
\clearpage
{\sffamily\color{BookAccent}\Large\bfseries Figure 3.9}\par\medskip
The unitary transforms $U^\alpha f$ for $0\leq\alpha\leq1$, using the fixed phases in the text. Real and imaginary parts are shown separately.\par\smallskip
{\small Book page 44; original figure 3.9.}\par\bigskip
\begin{minipage}[t]{.485\linewidth}\centering{\sffamily\bfseries Previous edition}\par\medskip
\includegraphics[width=\linewidth,height=.66\textheight,keepaspectratio]{comparator-figures/before/fig-root-interm-tfd.pdf}
\end{minipage}\hfill\begin{minipage}[t]{.485\linewidth}\centering{\sffamily\bfseries Revised edition}\par\medskip
\includegraphics[width=\linewidth,height=.66\textheight,keepaspectratio]{figures/chapter-discrete-fourier-transform/intermediate-fourier-transform/figure.pdf}
\end{minipage}
\clearpage
{\sffamily\color{BookAccent}\Large\bfseries Figure 3.10}\par\medskip
Periodic trigonometric interpolation of eleven samples. The faint polygon joins the samples for reference; it is not an underlying continuous signal.\par\smallskip
{\small Book page 46; original figure 3.10.}\par\bigskip
\begin{minipage}[t]{.485\linewidth}\centering{\sffamily\bfseries Previous edition}\par\medskip
\includegraphics[width=\linewidth,height=.66\textheight,keepaspectratio]{comparator-figures/before/fig-fft-interpolation.pdf}
\end{minipage}\hfill\begin{minipage}[t]{.485\linewidth}\centering{\sffamily\bfseries Revised edition}\par\medskip
\includegraphics[width=\linewidth,height=.66\textheight,keepaspectratio]{figures/chapter-discrete-fourier-transform/trigonometric-interpolation/figure.pdf}
\end{minipage}
\clearpage
{\sffamily\color{BookAccent}\Large\bfseries Figure 3.11}\par\medskip
The Chebyshev polynomials $T_k(x)=\cos(k\arccos x)$ for $k=1,\ldots,6$.\par\smallskip
{\small Book page 46; original figure 3.11.}\par\bigskip
\begin{minipage}[t]{.485\linewidth}\centering{\sffamily\bfseries Previous edition}\par\medskip
\includegraphics[width=\linewidth,height=.66\textheight,keepaspectratio]{comparator-figures/before/fig-polynomes-chebyshev.pdf}
\end{minipage}\hfill\begin{minipage}[t]{.485\linewidth}\centering{\sffamily\bfseries Revised edition}\par\medskip
\includegraphics[width=\linewidth,height=.66\textheight,keepaspectratio]{figures/chapter-discrete-fourier-transform/chebyshev-polynomials/figure.pdf}
\end{minipage}
\clearpage
{\sffamily\color{BookAccent}\Large\bfseries Figure 3.12}\par\medskip
Degree-ten polynomial interpolation of $f(x)=1/(0.3^2+x^2)$ at eleven equally spaced points and at the roots of $T_{11}$. Both are Lagrange interpolants; the choice of nodes controls the endpoint oscillations.\par\smallskip
{\small Book page 47; original figure 3.12.}\par\bigskip
\begin{minipage}[t]{.485\linewidth}\centering{\sffamily\bfseries Previous edition}\par\medskip
\includegraphics[width=\linewidth,height=.66\textheight,keepaspectratio]{comparator-figures/before/fig-interpolation-chebyshev.pdf}
\end{minipage}\hfill\begin{minipage}[t]{.485\linewidth}\centering{\sffamily\bfseries Revised edition}\par\medskip
\includegraphics[width=\linewidth,height=.66\textheight,keepaspectratio]{figures/chapter-discrete-fourier-transform/chebyshev-interpolation/figure.pdf}
\end{minipage}
\clearpage
{\sffamily\color{BookAccent}\Large\bfseries Figure 3.13}\par\medskip
Fractional derivatives of $e^{-x^2/2}$ for orders $0\leq\alpha\leq2$, computed from the multiplier $(\imath\xi)^\alpha$ with the phase convention in the text.\par\smallskip
{\small Book page 48; original figure 3.13.}\par\bigskip
\begin{minipage}[t]{.485\linewidth}\centering{\sffamily\bfseries Previous edition}\par\medskip
\includegraphics[width=\linewidth,height=.66\textheight,keepaspectratio]{comparator-figures/before/fig-derivation-fractionionnaire.pdf}
\end{minipage}\hfill\begin{minipage}[t]{.485\linewidth}\centering{\sffamily\bfseries Revised edition}\par\medskip
\includegraphics[width=\linewidth,height=.66\textheight,keepaspectratio]{figures/chapter-discrete-fourier-transform/fractional-differentiation/figure.pdf}
\end{minipage}
\clearpage
{\sffamily\color{BookAccent}\Large\bfseries Figure 3.14}\par\medskip
Entrywise magnitudes of the unitary matrices $U^\alpha$, on one common nonlinear color scale. At $\alpha=1$ every magnitude equals $1/\sqrt N$; at $\alpha=2$ the matrix is cyclic reversal.\par\smallskip
{\small Book page 48; original figure 3.14.}\par\bigskip
\begin{minipage}[t]{.485\linewidth}\centering{\sffamily\bfseries Previous edition}\par\medskip
\includegraphics[width=\linewidth,height=.66\textheight,keepaspectratio]{comparator-figures/before/fig-matrix-partial-tfd.pdf}
\end{minipage}\hfill\begin{minipage}[t]{.485\linewidth}\centering{\sffamily\bfseries Revised edition}\par\medskip
\includegraphics[width=\linewidth,height=.66\textheight,keepaspectratio]{figures/chapter-discrete-fourier-transform/fractional-fourier-matrices/figure.pdf}
\end{minipage}
\clearpage
{\sffamily\color{BookAccent}\Large\bfseries Figure 3.15}\par\medskip
Unitary fractional Fourier transforms of a discrete Gaussian. The fixed eigenspace phases give a continuous family; at $\alpha=2$ this even input is recovered.\par\smallskip
{\small Book page 49; original figure 3.15.}\par\bigskip
\begin{minipage}[t]{.485\linewidth}\centering{\sffamily\bfseries Previous edition}\par\medskip
\includegraphics[width=\linewidth,height=.66\textheight,keepaspectratio]{comparator-figures/before/fig-transforme-partial-fourier.pdf}
\end{minipage}\hfill\begin{minipage}[t]{.485\linewidth}\centering{\sffamily\bfseries Revised edition}\par\medskip
\includegraphics[width=\linewidth,height=.66\textheight,keepaspectratio]{figures/chapter-discrete-fourier-transform/fractional-fourier-gaussian/figure.pdf}
\end{minipage}
\clearpage
{\sffamily\color{BookAccent}\Large\bfseries Figure 3.16}\par\medskip
Magnitudes of the first three vectors in the structured DFT eigenbasis for $N=128$, ordered by decreasing eigenvalue of $S$. The index is centered for display.\par\smallskip
{\small Book page 50; original figure 3.16.}\par\bigskip
\begin{minipage}[t]{.485\linewidth}\centering{\sffamily\bfseries Previous edition}\par\medskip
\includegraphics[width=\linewidth,height=.66\textheight,keepaspectratio]{comparator-figures/before/fig-eigenvectors-tfd.pdf}
\end{minipage}\hfill\begin{minipage}[t]{.485\linewidth}\centering{\sffamily\bfseries Revised edition}\par\medskip
\includegraphics[width=\linewidth,height=.66\textheight,keepaspectratio]{figures/chapter-discrete-fourier-transform/fourier-eigenvectors/figure.pdf}
\end{minipage}
\clearpage
{\sffamily\color{BookAccent}\Large\bfseries Figure 3.17}\par\medskip
Magnitudes of the structured orthonormal eigenbasis for $N=128$. Columns follow decreasing eigenvalues of $S$, with even parity first in a tie; rows show centered sample indices.\par\smallskip
{\small Book page 50; original figure 3.17.}\par\bigskip
\begin{minipage}[t]{.485\linewidth}\centering{\sffamily\bfseries Previous edition}\par\medskip
\includegraphics[width=\linewidth,height=.66\textheight,keepaspectratio]{comparator-figures/before/fig-matrix-eigenvectors-tfd.pdf}
\end{minipage}\hfill\begin{minipage}[t]{.485\linewidth}\centering{\sffamily\bfseries Revised edition}\par\medskip
\includegraphics[width=\linewidth,height=.66\textheight,keepaspectratio]{figures/chapter-discrete-fourier-transform/fourier-eigenbasis-matrix/figure.pdf}
\end{minipage}
\clearpage
{\sffamily\color{BookAccent}\Large\bfseries Figure 3.18}\par\medskip
Orthogonalizing the cyclic translates of $f=\delta_0+a\delta_1$ on $\ZZ/20\ZZ$, for $a=0.5$ and $a=0.9$. The output $g$ satisfies $|\widehat g|=\sqrt{20}$, as required by the normalized inner product.\par\smallskip
{\small Book page 51; original figure 3.18.}\par\bigskip
\begin{minipage}[t]{.485\linewidth}\centering{\sffamily\bfseries Previous edition}\par\medskip
\includegraphics[width=\linewidth,height=.66\textheight,keepaspectratio]{comparator-figures/before/fig-orthogonalisation-fourier.pdf}
\end{minipage}\hfill\begin{minipage}[t]{.485\linewidth}\centering{\sffamily\bfseries Revised edition}\par\medskip
\includegraphics[width=\linewidth,height=.66\textheight,keepaspectratio]{figures/chapter-discrete-fourier-transform/orthogonalizing-translates/figure.pdf}
\end{minipage}
\clearpage
{\sffamily\color{BookAccent}\Large\bfseries Figure 4.1}\par\medskip
Sampling and zero padding for $f=\mathbf1_{[1/2,1)}$. Rows change the number $N$ of measured samples; columns change only the padded length $M$. The DFT magnitudes include the quadrature factor $1/N$; the faint curve is the continuous-transform magnitude.\par\smallskip
{\small Book page 54; original figure 4.1.}\par\bigskip
\begin{minipage}[t]{.485\linewidth}\centering{\sffamily\bfseries Previous edition}\par\medskip
\includegraphics[width=\linewidth,height=.66\textheight,keepaspectratio]{comparator-figures/before/fig-calcul-approach-tf.pdf}
\end{minipage}\hfill\begin{minipage}[t]{.485\linewidth}\centering{\sffamily\bfseries Revised edition}\par\medskip
\includegraphics[width=\linewidth,height=.66\textheight,keepaspectratio]{figures/chapter-fourier-applications/sampling-and-padding/figure.pdf}
\end{minipage}
\clearpage
{\sffamily\color{BookAccent}\Large\bfseries Figure 4.2}\par\medskip
Gaussian kernels and their action on one fixed irregular signal. Widths are in sample units, and the signal is reflected at its endpoints. The first column is the exact identity kernel.\par\smallskip
{\small Book page 55; original figure 4.2.}\par\bigskip
\begin{minipage}[t]{.485\linewidth}\centering{\sffamily\bfseries Previous edition}\par\medskip
\includegraphics[width=\linewidth,height=.66\textheight,keepaspectratio]{comparator-figures/before/fig-filter-smoothing.pdf}
\end{minipage}\hfill\begin{minipage}[t]{.485\linewidth}\centering{\sffamily\bfseries Revised edition}\par\medskip
\includegraphics[width=\linewidth,height=.66\textheight,keepaspectratio]{figures/chapter-fourier-applications/gaussian-signal-smoothing/figure.pdf}
\end{minipage}
\clearpage
{\sffamily\color{BookAccent}\Large\bfseries Figure 4.3}\par\medskip
Impulse, frequency, and step responses of three unit-sum symmetric FIR kernels. The frequency plots show a neighborhood of zero; the step responses are cumulative sums of the impulse responses.\par\smallskip
{\small Book page 55; original figure 4.3.}\par\bigskip
\begin{minipage}[t]{.485\linewidth}\centering{\sffamily\bfseries Previous edition}\par\medskip
\includegraphics[width=\linewidth,height=.66\textheight,keepaspectratio]{comparator-figures/before/fig-diff-types-responses.pdf}
\end{minipage}\hfill\begin{minipage}[t]{.485\linewidth}\centering{\sffamily\bfseries Revised edition}\par\medskip
\includegraphics[width=\linewidth,height=.66\textheight,keepaspectratio]{figures/chapter-fourier-applications/filter-responses/figure.pdf}
\end{minipage}
\clearpage
{\sffamily\color{BookAccent}\Large\bfseries Figure 4.4}\par\medskip
Discrete Gaussian kernels (top) and smoothing of the flower image (bottom), with reflected boundaries. Each kernel is displayed relative to its own peak; all image panels share the grayscale range $[0,1]$.\par\smallskip
{\small Book page 56; original figure 4.4.}\par\bigskip
\begin{minipage}[t]{.485\linewidth}\centering{\sffamily\bfseries Previous edition}\par\medskip
\includegraphics[width=\linewidth,height=.66\textheight,keepaspectratio]{comparator-figures/before/fig-Gaussian-filter-2.pdf}
\end{minipage}\hfill\begin{minipage}[t]{.485\linewidth}\centering{\sffamily\bfseries Revised edition}\par\medskip
\includegraphics[width=\linewidth,height=.66\textheight,keepaspectratio]{figures/chapter-fourier-applications/gaussian-image-smoothing/figure.pdf}
\end{minipage}
\clearpage
{\sffamily\color{BookAccent}\Large\bfseries Figure 4.5}\par\medskip
A flower image corrupted by additive Gaussian noise of standard deviation $0.12$, and Gaussian smoothing with width $0.9$ pixels. The displayed PSNR values use the same clean reference; intensity clipping is for display only.\par\smallskip
{\small Book page 57; original figure 4.5.}\par\bigskip
\begin{minipage}[t]{.485\linewidth}\centering{\sffamily\bfseries Previous edition}\par\medskip
\includegraphics[width=\linewidth,height=.66\textheight,keepaspectratio]{comparator-figures/before/fig-filter-2d.pdf}
\end{minipage}\hfill\begin{minipage}[t]{.485\linewidth}\centering{\sffamily\bfseries Revised edition}\par\medskip
\includegraphics[width=\linewidth,height=.66\textheight,keepaspectratio]{figures/chapter-fourier-applications/image-denoising/figure.pdf}
\end{minipage}
\clearpage
{\sffamily\color{BookAccent}\Large\bfseries Figure 4.6}\par\medskip
Four regimes of polygon filtering. Curves progress from the initial polygon to later iterates. Expanding active multipliers cause divergence; unit-modulus phases can keep the trajectory bounded without convergence.\par\smallskip
{\small Book page 58; original figure 4.6.}\par\bigskip
\begin{minipage}[t]{.485\linewidth}\centering{\sffamily\bfseries Previous edition}\par\medskip
\includegraphics[width=\linewidth,height=.66\textheight,keepaspectratio]{comparator-figures/before/fig-filter-polygons-1.pdf}
\end{minipage}\hfill\begin{minipage}[t]{.485\linewidth}\centering{\sffamily\bfseries Revised edition}\par\medskip
\includegraphics[width=\linewidth,height=.66\textheight,keepaspectratio]{figures/chapter-fourier-applications/polygon-iteration-regimes/figure.pdf}
\end{minipage}
\clearpage
{\sffamily\color{BookAccent}\Large\bfseries Figure 4.7}\par\medskip
Successive midpoint averaging converges to the vertex centroid (left). Keeping only frequency one produces a regular polygon of radius $|\widehat\Pi[1]|/N$ (right). Red marks indicate the limits.\par\smallskip
{\small Book page 59; original figure 4.7.}\par\bigskip
\begin{minipage}[t]{.485\linewidth}\centering{\sffamily\bfseries Previous edition}\par\medskip
\includegraphics[width=\linewidth,height=.66\textheight,keepaspectratio]{comparator-figures/before/fig-filter-polygons-2.pdf}
\end{minipage}\hfill\begin{minipage}[t]{.485\linewidth}\centering{\sffamily\bfseries Revised edition}\par\medskip
\includegraphics[width=\linewidth,height=.66\textheight,keepaspectratio]{figures/chapter-fourier-applications/polygon-filtering/figure.pdf}
\end{minipage}
\clearpage
{\sffamily\color{BookAccent}\Large\bfseries Figure 4.8}\par\medskip
A polygon and a perturbed, rotated version, with their magnitude-based Fourier descriptors. The perturbation changes the descriptor; the rotation alone would leave it unchanged.\par\smallskip
{\small Book page 60; original figure 4.8.}\par\bigskip
\begin{minipage}[t]{.485\linewidth}\centering{\sffamily\bfseries Previous edition}\par\medskip
\includegraphics[width=\linewidth,height=.66\textheight,keepaspectratio]{comparator-figures/before/fig-fourier-descriptors.pdf}
\end{minipage}\hfill\begin{minipage}[t]{.485\linewidth}\centering{\sffamily\bfseries Revised edition}\par\medskip
\includegraphics[width=\linewidth,height=.66\textheight,keepaspectratio]{figures/chapter-fourier-applications/fourier-descriptors/figure.pdf}
\end{minipage}
\clearpage
{\sffamily\color{BookAccent}\Large\bfseries Figure 4.9}\par\medskip
Spectral heat evolution of an interval indicator on the unit circle, with conductivity $\kappa=1/2$. The mean is preserved and the nonconstant Fourier modes decay.\par\smallskip
{\small Book page 61; original figure 4.9.}\par\bigskip
\begin{minipage}[t]{.485\linewidth}\centering{\sffamily\bfseries Previous edition}\par\medskip
\includegraphics[width=\linewidth,height=.66\textheight,keepaspectratio]{comparator-figures/before/fig-eq-heat-evolution.pdf}
\end{minipage}\hfill\begin{minipage}[t]{.485\linewidth}\centering{\sffamily\bfseries Revised edition}\par\medskip
\includegraphics[width=\linewidth,height=.66\textheight,keepaspectratio]{figures/chapter-fourier-applications/heat-evolution/figure.pdf}
\end{minipage}
\clearpage
{\sffamily\color{BookAccent}\Large\bfseries Figure 4.10}\par\medskip
Dirichlet Poisson solution with reference $u(x,y)=e^{xy}$, on a grid with $N=32$ intervals per axis.\par\smallskip
{\small Book page 62; original figure 4.10.}\par\bigskip
\begin{minipage}[t]{.485\linewidth}\centering{\sffamily\bfseries Previous edition}\par\medskip
\includegraphics[width=\linewidth,height=.66\textheight,keepaspectratio]{comparator-figures/before/fig-eq-fish.pdf}
\end{minipage}\hfill\begin{minipage}[t]{.485\linewidth}\centering{\sffamily\bfseries Revised edition}\par\medskip
\includegraphics[width=\linewidth,height=.66\textheight,keepaspectratio]{figures/chapter-fourier-applications/poisson-reference/figure.pdf}
\end{minipage}
\clearpage
{\sffamily\color{BookAccent}\Large\bfseries Figure 4.11}\par\medskip
The same source with boundary values lowered by $0.45\sin(\pi s)$ on $x=0$ and $y=0$. All four corner values remain compatible.\par\smallskip
{\small Book page 62; original figure 4.11.}\par\bigskip
\begin{minipage}[t]{.485\linewidth}\centering{\sffamily\bfseries Previous edition}\par\medskip
\includegraphics[width=\linewidth,height=.66\textheight,keepaspectratio]{comparator-figures/before/fig-eq-Poisson-modif.pdf}
\end{minipage}\hfill\begin{minipage}[t]{.485\linewidth}\centering{\sffamily\bfseries Revised edition}\par\medskip
\includegraphics[width=\linewidth,height=.66\textheight,keepaspectratio]{figures/chapter-fourier-applications/poisson-modified-boundary/figure.pdf}
\end{minipage}
\clearpage
{\sffamily\color{BookAccent}\Large\bfseries Figure 4.12}\par\medskip
Coefficient vectors of $P$, $Q$, and their product. Length-eleven FFT convolution agrees with the ordinary degree-ten polynomial product.\par\smallskip
{\small Book page 63; original figure 4.12.}\par\bigskip
\begin{minipage}[t]{.485\linewidth}\centering{\sffamily\bfseries Previous edition}\par\medskip
\includegraphics[width=\linewidth,height=.66\textheight,keepaspectratio]{comparator-figures/before/fig-convolution-polynomial.pdf}
\end{minipage}\hfill\begin{minipage}[t]{.485\linewidth}\centering{\sffamily\bfseries Revised edition}\par\medskip
\includegraphics[width=\linewidth,height=.66\textheight,keepaspectratio]{figures/chapter-fourier-applications/polynomial-convolution/figure.pdf}
\end{minipage}
\clearpage
{\sffamily\color{BookAccent}\Large\bfseries Figure 4.13}\par\medskip
Three time steps for the same flower image, with $N=128$ and $h=0.0002$. Explicit Euler is unstable; Crank--Nicolson is stable but weakly damps high frequencies; backward Euler damps them strongly. All images use $[0,1]$, and each panel reports its actual value range.\par\smallskip
{\small Book page 65; original figure 4.13.}\par\bigskip
\begin{minipage}[t]{.485\linewidth}\centering{\sffamily\bfseries Previous edition}\par\medskip
\includegraphics[width=\linewidth,height=.66\textheight,keepaspectratio]{comparator-figures/before/fig-eq-heat-finite-diff.pdf}
\end{minipage}\hfill\begin{minipage}[t]{.485\linewidth}\centering{\sffamily\bfseries Revised edition}\par\medskip
\includegraphics[width=\linewidth,height=.66\textheight,keepaspectratio]{figures/chapter-fourier-applications/heat-time-stepping/figure.pdf}
\end{minipage}
\clearpage
{\sffamily\color{BookAccent}\Large\bfseries Figure 4.14}\par\medskip
Template matching in the flower image. The red box identifies the extracted patch; crosses mark its location in the raw and normalized correlation maps. Only complete-template positions are included.\par\smallskip
{\small Book page 67; original figure 4.14.}\par\bigskip
\begin{minipage}[t]{.485\linewidth}\centering{\sffamily\bfseries Previous edition}\par\medskip
\includegraphics[width=\linewidth,height=.66\textheight,keepaspectratio]{comparator-figures/before/fig-correlation-2d.pdf}
\end{minipage}\hfill\begin{minipage}[t]{.485\linewidth}\centering{\sffamily\bfseries Revised edition}\par\medskip
\includegraphics[width=\linewidth,height=.66\textheight,keepaspectratio]{figures/chapter-fourier-applications/template-matching/figure.pdf}
\end{minipage}
\clearpage
{\sffamily\color{BookAccent}\Large\bfseries Figure 4.15}\par\medskip
Rotations of the flower image by three Fourier shears. The odd grid avoids a Nyquist ambiguity; a white margin limits wraparound. Cropping and intensity clipping are applied only to the displayed results.\par\smallskip
{\small Book page 68; original figure 4.15.}\par\bigskip
\begin{minipage}[t]{.485\linewidth}\centering{\sffamily\bfseries Previous edition}\par\medskip
\includegraphics[width=\linewidth,height=.66\textheight,keepaspectratio]{comparator-figures/before/fig-fft-rotation.pdf}
\end{minipage}\hfill\begin{minipage}[t]{.485\linewidth}\centering{\sffamily\bfseries Revised edition}\par\medskip
\includegraphics[width=\linewidth,height=.66\textheight,keepaspectratio]{figures/chapter-fourier-applications/fourier-image-rotation/figure.pdf}
\end{minipage}
\clearpage
{\sffamily\color{BookAccent}\Large\bfseries Figure 4.16}\par\medskip
Low-pass filters with increasingly wide raised-cosine transitions. The columns show the centered impulse response, the prescribed DFT samples, and the densely sampled FIR frequency response.\par\smallskip
{\small Book page 68; original figure 4.16.}\par\bigskip
\begin{minipage}[t]{.485\linewidth}\centering{\sffamily\bfseries Previous edition}\par\medskip
\includegraphics[width=\linewidth,height=.66\textheight,keepaspectratio]{comparator-figures/before/fig-low-pass-filter.pdf}
\end{minipage}\hfill\begin{minipage}[t]{.485\linewidth}\centering{\sffamily\bfseries Revised edition}\par\medskip
\includegraphics[width=\linewidth,height=.66\textheight,keepaspectratio]{figures/chapter-fourier-applications/low-pass-filters/figure.pdf}
\end{minipage}
\clearpage
{\sffamily\color{BookAccent}\Large\bfseries Figure 4.17}\par\medskip
Centered B-splines of degrees zero through three. Their support is $[-(n+1)/2,(n+1)/2]$.\par\smallskip
{\small Book page 69; original figure 4.17.}\par\bigskip
\begin{minipage}[t]{.485\linewidth}\centering{\sffamily\bfseries Previous edition}\par\medskip
\includegraphics[width=\linewidth,height=.66\textheight,keepaspectratio]{comparator-figures/before/fig-box-spline.pdf}
\end{minipage}\hfill\begin{minipage}[t]{.485\linewidth}\centering{\sffamily\bfseries Revised edition}\par\medskip
\includegraphics[width=\linewidth,height=.66\textheight,keepaspectratio]{figures/chapter-fourier-applications/centered-bsplines/figure.pdf}
\end{minipage}
\clearpage
{\sffamily\color{BookAccent}\Large\bfseries Figure 4.18}\par\medskip
The cubic cardinal kernel, the cubic B-spline, and sinc. The cardinal kernel and sinc equal one at zero and vanish at every other integer; the cubic B-spline requires a prefilter.\par\smallskip
{\small Book page 70; original figure 4.18.}\par\bigskip
\begin{minipage}[t]{.485\linewidth}\centering{\sffamily\bfseries Previous edition}\par\medskip
\includegraphics[width=\linewidth,height=.66\textheight,keepaspectratio]{comparator-figures/before/fig-comparison-spline-sinc.pdf}
\end{minipage}\hfill\begin{minipage}[t]{.485\linewidth}\centering{\sffamily\bfseries Revised edition}\par\medskip
\includegraphics[width=\linewidth,height=.66\textheight,keepaspectratio]{figures/chapter-fourier-applications/cardinal-spline-and-sinc/figure.pdf}
\end{minipage}
\clearpage
{\sffamily\color{BookAccent}\Large\bfseries Figure 4.19}\par\medskip
Cubic interpolation of the same eight samples with natural and not-a-knot boundary conditions. Natural conditions set the endpoint second derivatives to zero.\par\smallskip
{\small Book page 70; original figure 4.19.}\par\bigskip
\begin{minipage}[t]{.485\linewidth}\centering{\sffamily\bfseries Previous edition}\par\medskip
\includegraphics[width=\linewidth,height=.66\textheight,keepaspectratio]{comparator-figures/before/fig-spline-cubic.pdf}
\end{minipage}\hfill\begin{minipage}[t]{.485\linewidth}\centering{\sffamily\bfseries Revised edition}\par\medskip
\includegraphics[width=\linewidth,height=.66\textheight,keepaspectratio]{figures/chapter-fourier-applications/cubic-spline-boundaries/figure.pdf}
\end{minipage}
\clearpage
{\sffamily\color{BookAccent}\Large\bfseries Figure 5.1}\par\medskip
A triangular signal and its densely evaluated Fourier and Hartley spectra. The Hartley spectrum equals the real part minus the imaginary part of the DFT. Zero padding changes only the frequency grid.\par\smallskip
{\small Book page 72; original figure 5.1.}\par\bigskip
\begin{minipage}[t]{.485\linewidth}\centering{\sffamily\bfseries Previous edition}\par\medskip
\includegraphics[width=\linewidth,height=.66\textheight,keepaspectratio]{comparator-figures/before/fig-comparison-fht-fft.pdf}
\end{minipage}\hfill\begin{minipage}[t]{.485\linewidth}\centering{\sffamily\bfseries Revised edition}\par\medskip
\includegraphics[width=\linewidth,height=.66\textheight,keepaspectratio]{figures/chapter-related-transforms/hartley-and-fourier/figure.pdf}
\end{minipage}
\clearpage
{\sffamily\color{BookAccent}\Large\bfseries Figure 5.2}\par\medskip
Zeros on the unit circle and poles at radius $0.9$ for a notch filter at angular frequencies $\pm\pi/4$. Circles mark zeros; crosses mark poles.\par\smallskip
{\small Book page 75; original figure 5.2.}\par\bigskip
\begin{minipage}[t]{.485\linewidth}\centering{\sffamily\bfseries Previous edition}\par\medskip
\includegraphics[width=\linewidth,height=.66\textheight,keepaspectratio]{comparator-figures/before/fig-location-poles-notch-filter.pdf}
\end{minipage}\hfill\begin{minipage}[t]{.485\linewidth}\centering{\sffamily\bfseries Revised edition}\par\medskip
\includegraphics[width=\linewidth,height=.66\textheight,keepaspectratio]{figures/chapter-related-transforms/notch-poles-and-zeros/figure.pdf}
\end{minipage}
\clearpage
{\sffamily\color{BookAccent}\Large\bfseries Figure 5.3}\par\medskip
Magnitude of the rational notch-filter transfer function in the complex plane. Values above two are clipped to reveal the zeros and pole neighborhoods. The causal Laurent series converges for $|z|>0.9$.\par\smallskip
{\small Book page 75; original figure 5.3.}\par\bigskip
\begin{minipage}[t]{.485\linewidth}\centering{\sffamily\bfseries Previous edition}\par\medskip
\includegraphics[width=\linewidth,height=.66\textheight,keepaspectratio]{comparator-figures/before/fig-notch-filter-trans-z.pdf}
\end{minipage}\hfill\begin{minipage}[t]{.485\linewidth}\centering{\sffamily\bfseries Revised edition}\par\medskip
\includegraphics[width=\linewidth,height=.66\textheight,keepaspectratio]{figures/chapter-related-transforms/notch-transfer-surface/figure.pdf}
\end{minipage}
\clearpage
{\sffamily\color{BookAccent}\Large\bfseries Figure 5.4}\par\medskip
Frequency-response magnitude and causal impulse response of the notch filter. The response vanishes at angular frequency $\pi/4$.\par\smallskip
{\small Book page 76; original figure 5.4.}\par\bigskip
\begin{minipage}[t]{.485\linewidth}\centering{\sffamily\bfseries Previous edition}\par\medskip
\includegraphics[width=\linewidth,height=.66\textheight,keepaspectratio]{comparator-figures/before/fig-notch-filter-rep-imp-freq.pdf}
\end{minipage}\hfill\begin{minipage}[t]{.485\linewidth}\centering{\sffamily\bfseries Revised edition}\par\medskip
\includegraphics[width=\linewidth,height=.66\textheight,keepaspectratio]{figures/chapter-related-transforms/notch-responses/figure.pdf}
\end{minipage}
\clearpage
{\sffamily\color{BookAccent}\Large\bfseries Figure 5.5}\par\medskip
A series RLC circuit, with the output voltage $V_s$ measured across the resistor and the input $V_e$ across the two left terminals.\par\smallskip
{\small Book page 76; original figure 5.5.}\par\bigskip
\begin{minipage}[t]{.485\linewidth}\centering{\sffamily\bfseries Previous edition}\par\medskip
\includegraphics[width=\linewidth,height=.66\textheight,keepaspectratio]{comparator-figures/before/fig-circuit-rlc.pdf}
\end{minipage}\hfill\begin{minipage}[t]{.485\linewidth}\centering{\sffamily\bfseries Revised edition}\par\medskip
\includegraphics[width=\linewidth,height=.66\textheight,keepaspectratio]{figures/chapter-related-transforms/rlc-circuit/figure.pdf}
\end{minipage}
\clearpage
{\sffamily\color{BookAccent}\Large\bfseries Figure 5.6}\par\medskip
The notch-filter transfer function evaluated on a unit circle and an outward logarithmic spiral. Both contours lie in the causal region of convergence; leaving the unit circle changes the depth of the notches.\par\smallskip
{\small Book page 78; original figure 5.6.}\par\bigskip
\begin{minipage}[t]{.485\linewidth}\centering{\sffamily\bfseries Previous edition}\par\medskip
\includegraphics[width=\linewidth,height=.66\textheight,keepaspectratio]{comparator-figures/before/fig-z-transform.pdf}
\end{minipage}\hfill\begin{minipage}[t]{.485\linewidth}\centering{\sffamily\bfseries Revised edition}\par\medskip
\includegraphics[width=\linewidth,height=.66\textheight,keepaspectratio]{figures/chapter-related-transforms/z-transform-contours/figure.pdf}
\end{minipage}
\clearpage
{\sffamily\color{BookAccent}\Large\bfseries Figure 5.7}\par\medskip
Scaled-frequency Fourier transforms of an interval indicator. The real and imaginary parts share the same scale. Changing $\alpha$ changes the sampled frequency grid.\par\smallskip
{\small Book page 79; original figure 5.7.}\par\bigskip
\begin{minipage}[t]{.485\linewidth}\centering{\sffamily\bfseries Previous edition}\par\medskip
\includegraphics[width=\linewidth,height=.66\textheight,keepaspectratio]{comparator-figures/before/fig-transforme-fourier-fractionionnaire-1d.pdf}
\end{minipage}\hfill\begin{minipage}[t]{.485\linewidth}\centering{\sffamily\bfseries Revised edition}\par\medskip
\includegraphics[width=\linewidth,height=.66\textheight,keepaspectratio]{figures/chapter-related-transforms/scaled-fourier-signal/figure.pdf}
\end{minipage}
\clearpage
{\sffamily\color{BookAccent}\Large\bfseries Figure 5.8}\par\medskip
Finite-record spectra of sinusoids with integer and noninteger numbers of cycles. The continuous curves evaluate the finite sum; the dots are DFT samples. The vertical line marks the generating frequency.\par\smallskip
{\small Book page 79; original figure 5.8.}\par\bigskip
\begin{minipage}[t]{.485\linewidth}\centering{\sffamily\bfseries Previous edition}\par\medskip
\includegraphics[width=\linewidth,height=.66\textheight,keepaspectratio]{comparator-figures/before/fig-spectrum-non-periodic-signal.pdf}
\end{minipage}\hfill\begin{minipage}[t]{.485\linewidth}\centering{\sffamily\bfseries Revised edition}\par\medskip
\includegraphics[width=\linewidth,height=.66\textheight,keepaspectratio]{figures/chapter-related-transforms/spectral-leakage/figure.pdf}
\end{minipage}
\clearpage
{\sffamily\color{BookAccent}\Large\bfseries Figure 5.9}\par\medskip
A noisy cosine, its coarse spectrum, a local spectral zoom, and an adjusted spectrum after shortening the record using the estimated peak. Spectral magnitudes are divided by the respective record lengths. The generating frequency is distinguished from its estimate.\par\smallskip
{\small Book page 80; original figure 5.9.}\par\bigskip
\begin{minipage}[t]{.485\linewidth}\centering{\sffamily\bfseries Previous edition}\par\medskip
\includegraphics[width=\linewidth,height=.66\textheight,keepaspectratio]{comparator-figures/before/fig-transforme-fourier-fractionionnaire.pdf}
\end{minipage}\hfill\begin{minipage}[t]{.485\linewidth}\centering{\sffamily\bfseries Revised edition}\par\medskip
\includegraphics[width=\linewidth,height=.66\textheight,keepaspectratio]{figures/chapter-related-transforms/spectral-zoom-and-adjustment/figure.pdf}
\end{minipage}
\clearpage
{\sffamily\color{BookAccent}\Large\bfseries Figure 5.10}\par\medskip
Real and imaginary parts of fixed-branch unitary Hartley powers $V^t f$, with $V=\Hh/\sqrt N$. All panels use the same amplitude scale and preserve the norm.\par\smallskip
{\small Book page 81; original figure 5.10.}\par\bigskip
\begin{minipage}[t]{.485\linewidth}\centering{\sffamily\bfseries Previous edition}\par\medskip
\includegraphics[width=\linewidth,height=.66\textheight,keepaspectratio]{comparator-figures/before/fig-transfo-hartley-interm.pdf}
\end{minipage}\hfill\begin{minipage}[t]{.485\linewidth}\centering{\sffamily\bfseries Revised edition}\par\medskip
\includegraphics[width=\linewidth,height=.66\textheight,keepaspectratio]{figures/chapter-related-transforms/fractional-hartley/figure.pdf}
\end{minipage}
\clearpage
{\sffamily\color{BookAccent}\Large\bfseries Figure 5.11}\par\medskip
Generalized Hartley transforms of the same triangular signal. The phase $\lambda$ changes the relative cosine and sine contributions; the two endpoint panels illustrate the $2\pi$ periodicity.\par\smallskip
{\small Book page 81; original figure 5.11.}\par\bigskip
\begin{minipage}[t]{.485\linewidth}\centering{\sffamily\bfseries Previous edition}\par\medskip
\includegraphics[width=\linewidth,height=.66\textheight,keepaspectratio]{comparator-figures/before/fig-transforme-hartley-generalized.pdf}
\end{minipage}\hfill\begin{minipage}[t]{.485\linewidth}\centering{\sffamily\bfseries Revised edition}\par\medskip
\includegraphics[width=\linewidth,height=.66\textheight,keepaspectratio]{figures/chapter-related-transforms/generalized-hartley/figure.pdf}
\end{minipage}
\clearpage
{\sffamily\color{BookAccent}\Large\bfseries Figure 5.12}\par\medskip
Magnitudes of scaled-frequency two-dimensional Fourier transforms of a disk. Spectral panels share the contrast mapping $\log(1+100|G_2|/\sum f)/\log(101)$, shown dark on white; the input panel uses ordinary grayscale.\par\smallskip
{\small Book page 84; original figure 5.12.}\par\bigskip
\begin{minipage}[t]{.485\linewidth}\centering{\sffamily\bfseries Previous edition}\par\medskip
\includegraphics[width=\linewidth,height=.66\textheight,keepaspectratio]{comparator-figures/before/fig-transforme-fourier-fractionionnaire-2d.pdf}
\end{minipage}\hfill\begin{minipage}[t]{.485\linewidth}\centering{\sffamily\bfseries Revised edition}\par\medskip
\includegraphics[width=\linewidth,height=.66\textheight,keepaspectratio]{figures/chapter-related-transforms/scaled-fourier-image/figure.pdf}
\end{minipage}
\clearpage
{\sffamily\color{BookAccent}\Large\bfseries Figure 6.1}\par\medskip
Hamming distance on $\FF_2^3$. Each cube edge changes one coordinate; the highlighted shortest path has three edges.\par\smallskip
{\small Book page 89; original figure 6.1.}\par\bigskip
\begin{minipage}[t]{.485\linewidth}\centering{\sffamily\bfseries Previous edition}\par\medskip
\includegraphics[width=\linewidth,height=.66\textheight,keepaspectratio]{comparator-figures/before/fig-visual-dist-hamming.pdf}
\end{minipage}\hfill\begin{minipage}[t]{.485\linewidth}\centering{\sffamily\bfseries Revised edition}\par\medskip
\includegraphics[width=\linewidth,height=.66\textheight,keepaspectratio]{figures/chapter-finite-fields-and-codes/hamming-cube/figure.pdf}
\end{minipage}
\clearpage
{\sffamily\color{BookAccent}\Large\bfseries Figure 7.1}\par\medskip
Wavelets on $\FF_{53}$: the signal, the centered wavelet, and the magnitude of its affine-group transform. Translation $b$ is horizontal and nonzero dilation $a$ is vertical. Each row has its own indicated magnitude scale. The first signal has nonzero mean, which the mean-zero wavelet removes.\par\smallskip
{\small Book page 110; original figure 7.1.}\par\bigskip
\begin{minipage}[t]{.485\linewidth}\centering{\sffamily\bfseries Previous edition}\par\medskip
\includegraphics[width=\linewidth,height=.66\textheight,keepaspectratio]{comparator-figures/before/wavelets-finite-body.pdf}
\end{minipage}\hfill\begin{minipage}[t]{.485\linewidth}\centering{\sffamily\bfseries Revised edition}\par\medskip
\includegraphics[width=\linewidth,height=.66\textheight,keepaspectratio]{figures/chapter-linear-representations/finite-field-wavelets/figure.pdf}
\end{minipage}
\clearpage
{\sffamily\color{BookAccent}\Large\bfseries Figure A.1}\par\medskip
Magnitudes of the rectangle, trapezoidal, and Simpson integration filters for $\Delta=1$. The poles at zero frequency and the additional Simpson pole at the Nyquist frequency are excluded from the plotted grid.\par\smallskip
{\small Book page 148; original figure A.1.}\par\bigskip
\begin{minipage}[t]{.485\linewidth}\centering{\sffamily\bfseries Previous edition}\par\medskip
\includegraphics[width=\linewidth,height=.66\textheight,keepaspectratio]{comparator-figures/before/fig-meth-quadratures-recursive-filter.pdf}
\end{minipage}\hfill\begin{minipage}[t]{.485\linewidth}\centering{\sffamily\bfseries Revised edition}\par\medskip
\includegraphics[width=\linewidth,height=.66\textheight,keepaspectratio]{figures/chapter-related-transforms/quadrature-filter-responses/figure.pdf}
\end{minipage}
\end{document}
